\section{Conclusion}

This project set out to undo TDA's calibration loss with three training-free, label-free
changes to its cache loop, and measured each one rather than assuming it worked. Margin-based
cache admission fails outright: its fixed logit threshold does not port across datasets, a
starved cache is a biased cache, and the top-two gap cannot tell a confident mistake from a
confident success, so the gate removes the wrong samples. We dropped it. Probabilistic fusion
fails as proposed, for a reason the equation makes visible once the two inputs' scales are
compared: a nearly flat cache distribution averaged into CLIP's peaked one caps every
prediction's confidence near one half, and ECE punishes that underconfidence exactly as it
punishes overconfidence. The repair follows from the diagnosis. Fusing only where CLIP is torn
between two classes, and leaving the cache out everywhere else, removes the dilution while
keeping the one thing the cache is good for, breaking ties. Gated fusion lowers TDA's mean ECE
from 11.24\% to 4.61\%, below zero-shot CLIP's own 5.57\%, at a mean accuracy cost of 1.62
points, most of it on EuroSAT. Leave-one-out temperature scaling helps on its own, cutting mean
ECE to 8.54\% at unchanged accuracy, but reverses on the dataset that was already calibrated
and adds nothing on top of the gate, because the gate has already removed the overconfidence it
was built to correct.

The method we recommend is therefore TDA's cache with gated probabilistic fusion: one extra
comparison per sample, no labels, no training, no change to the cache itself. Against the
project's own success criterion it passes on one dataset of five, and we report that count as
it stands. What the result supports is a narrower and, we think, more useful claim than the one
we set out to make: that TDA's calibration loss comes from consulting the cache on samples
where it can only add confidence, and that the gap between CLIP's top two scores, applied at
prediction time rather than at admission time, is enough to tell those samples apart.
